\section{Related Work}\label{sec:relatedwork}
Previous anomaly detection approaches are usually based on system logs~\cite{Deeplog, LogAnomaly, LogRobust, SwissLog, BertLog} or KPIs~\cite{Adsketch, SRCNN, InterFusion, USAD, OmniAnomaly}, or both~\cite{Hades}, targeting traditional distributed systems without complex invocation relationships.
Recently, some studies~\cite{TraceAnomaly, MMTrace, AID} have been presented to automate anomaly detection in microservice systems. 
\cite{TraceAnomaly} proposed to employ a variational autoencoder with a Bayes model to detect anomalies reflected by latency.
\cite{MMTrace} extracted operation sequence and invocation latency from traces and fed them into a multi-modal LSTM to identify anomalies.
These anomaly detection methods rely on single-source data (i.e., traces) and ignore other informative data such as logs and KPIs.

Tremendous efforts~\cite{TBAC, NetMedic, MonitorRank, CloudRanger, MicroCause, AutoMAP, TraceRCA} have been devoted to root cause localization in microservice or service-oriented systems, most of which rely on traces only and leverage traditional or naive machine learning techniques.
For example, \cite{MEPFL} conducted manual feature engineering in trace logs to predict latent errors and identify the faulty microservice via a decision tree. 
\cite{MicroHECL} proposed a high-efficient approach that dynamically constructs a service call graph and ranks candidate root causes based on correlation analysis. 
A recent study~\cite{DyCause} designed a crowd-sourcing solution to resolve user-space diagnosis for microservice kernel failures. 
These methods work well when the latent features of microservices are readily comprehensible but may lack scalability to larger-scale microservice systems with more complex features.
Deep learning-based approaches explore meaningful features from historical data to avoid manual feature engineering. Though deep learning has not been applied to root cause localization as far as we know, some approaches incorporated it for performance debugging. 
For example, to handle traces, \cite{Seer} used convolution networks and LSTM, and \cite{Sage} leveraged causal Bayesian networks.

However, they rely on traces and ignore other data sources, such as logs and KPIs, that can also reflect the microservice status. 
Also, they either focus on anomaly detection or root cause localization leading to the disconnection in the two closely related tasks. The inaccurate results of naive anomaly detectors affect the effectiveness of downstream localization.
Moreover, many methods combine manual feature engineering with traditional algorithms, making it insufficiently practical in large-scale systems.